% Generated draft: requires mathematical and source review.
\documentclass[11pt]{article}
\usepackage[margin=25mm]{geometry}
\usepackage{fontspec}
\setmainfont{Latin Modern Roman}
\usepackage{amsmath,amssymb,mathtools,unicode-math}
\setmathfont{Latin Modern Math}
\usepackage{enumitem,needspace,xurl,hyperref}
\setlength{\parindent}{0pt}
\setlength{\parskip}{5pt}
\setlength{\emergencystretch}{3em}
\newcommand{\N}{\mathbb{N}}
\newcommand{\Z}{\mathbb{Z}}
\newcommand{\Q}{\mathbb{Q}}
\newcommand{\R}{\mathbb{R}}
\newcommand{\C}{\mathbb{C}}
\newcommand{\F}{\mathbb{F}}
\newcommand{\A}{\mathbb{A}}
\newcommand{\PP}{\mathbb{P}}
\newcommand{\E}{\mathbb{E}}
\newcommand{\eps}{\varepsilon}
\newcommand{\dd}{\,\mathrm{d}}
\begin{document}
\section*{UnsolvedMath ID: 30006624}
\textbf{Source problem code:} catalogue entry
\subsection{Definitions and mathematical statement}
Let \((X,d)\) be a metric space. For \(a,b\in X\), its metric interval is
\[
I_d(a,b)=\{u\in X:d(a,b)=d(a,u)+d(u,b)\}.
\]
A median of a triple \((a,b,c)\in X^3\) is a point in
\[
I_d(a,b)\cap I_d(b,c)\cap I_d(c,a).
\]
The space is a \emph{median space} if this intersection contains exactly one point for every triple in \(X^3\). A path metric is a metric for which the distance between two points is the infimum of the lengths of continuous rectifiable paths joining them. Completeness means that every Cauchy sequence converges in \(X\). Simply connected means path connected with every loop null-homotopic.

For a precise reading of the supplied local condition, give each neighborhood the restricted metric \(d|_{U\times U}\). Under this reading, pointwise local medianity means
\[
\forall p\in X\ \exists U_p\subseteq X\text{ open},\quad p\in U_p,\qquad
\forall(a,b,c)\in U_p^3,\quad
\#\bigl(U_p\cap I_d(a,b)\cap I_d(b,c)\cap I_d(c,a)\bigr)=1.
\]
Here the median must lie in \(U_p\), and the distance equalities use the ambient metric \(d\).

\textbf{Question (supplied clean formulation).} Suppose that \((X,d)\) is a complete, simply connected path-metric space satisfying this pointwise local median condition. Must
\[
\forall(a,b,c)\in X^3,\qquad
\#\bigl(I_d(a,b)\cap I_d(b,c)\cap I_d(c,a)\bigr)=1?
\]
No common positive lower bound on the sizes of the local neighborhoods is assumed; the question includes spaces for which no such uniform bound is available.

\textbf{Unresolved qualification.} The restricted-metric interpretation above makes the clean formulation precise, but the supplied excerpt does not establish that it is the intended convention. Using the intrinsic path metric of a neighborhood, or allowing a local triple's median to lie outside that neighborhood, gives a different local condition. The original excerpt asks to remove the word ``uniform'' from an antecedent statement that is not supplied. Consequently, the exact original uniform hypothesis and its relationship to this formulation require source review. The original field also appends separate group-theoretic questions, recorded in the notes; they are not hypotheses of this metric-space question.
\subsection{Short English statement}
The supplied question asks whether a complete, simply connected length space is globally median whenever each point has a neighborhood in which every triple has exactly one median, without a uniform neighborhood size. The precise meaning of a median within a neighborhood and the omitted antecedent of the original question require review. The supplied open status is a dated curation claim.
\subsection{Sources}
\begin{itemize}
\item Supplied local UnsolvedMath registry record, permanent ID \(30006624\), title \emph{Cartan-Hadamard Theorem without Uniform Local Medianity}; exact source problem code \texttt{catalogue entry}. The record belongs to the Oberwolfach Reports Open Problems collection and gives proposed year \(2026\), with no named proposer in its proposer field.
\item Supplied report citation: \emph{Median Geometry and Applications} (2026), Oberwolfach Reports, pp.~487--540, DOI \texttt{10.4171/owr/2026/8}. Supplied source link: \url{https://doi.org/10.4171/owr/2026/8}. The background describes Oberwolfach Reports as documenting mathematical workshops and their open-problem discussions.
\item Supplied literature source, ``OWR report'': \url{https://ems.press/content/serial-article-files/53603}.
\item Supplied literature source, ``Bowditch preprint'': \url{https://bhbowditch.com/papers/ch-median.pdf}. The supplied assessment attributes the Cartan--Hadamard question to Bowditch; no preprint title, publication date, or more specific location is supplied.
\item Literature-check date: \texttt{2026-08-22}. The supplied dated curation labels the problem open and states that the cited preprint and report leave unresolved whether pointwise local uniqueness of medians suffices without a uniform radius. This is supplied dated curation, not an independently verified literature assessment. The background explicitly requests independent verification of time-sensitive status and cited claims.
\end{itemize}
\end{document}
